\section{Introduction}
\label{sec:intro}

Large language models answer questions through a single forward pass whose output depends on two design choices that the user controls: the prompt that specifies the task, and the decoding configuration that maps model probabilities into text. Both choices turn out to be fragile. A long line of work on prompt sensitivity shows that semantically equivalent reformulations of the same question can change model answers by large margins \citep{mizrahi2024state,sclar2024quantifying,leidinger2023language,zhu2023promptbench}. A parallel line of work on test-time compute shows that a single greedy decode is an unreliable summary of what the model knows, and that aggregating several decodes of the same prompt, as in self-consistency \citep{wang2023selfconsistency}, improves reasoning accuracy at the cost of additional inference calls.

These two observations motivate two different ways of spending a fixed budget of model calls. The first spends the budget on the answer space: keep the prompt fixed, sample several decodes at nonzero temperature, and aggregate them by majority vote. The second spends the budget on the prompt space: keep decoding deterministic or lightly stochastic, present several semantically equivalent paraphrases of the task, and aggregate the resulting answers. Both strategies are used in practice, both appear in the evaluation literature under separate headings, and both consume the same scarce resource. Surprisingly, the literature offers little guidance on which source of diversity yields more accuracy per call, whether the two sources detect the same errors, and how the answer changes with model capability and task difficulty.

The question is becoming more urgent as model behavior shifts. Post-trained flagship models are increasingly confident: on many legacy benchmarks they answer near ceiling, and preliminary reports suggest that repeated sampling at moderate temperature often returns identical answers, which would render answer-space ensembling inert precisely where deployments are most common. At the same time, uncertainty estimation methods built on sampling, such as semantic entropy \citep{kuhn2023semantic,farquhar2024detecting}, implicitly assume that the sample distribution carries information about correctness. If that assumption fails for confident models, practitioners need to know before they build on it.

This paper presents, to our knowledge, the first budget-matched, within-question comparison of answer-space and prompt-space ensembling across multiple model families and task types. Every question in our protocol is answered sixteen times by every model under matched conditions: one greedy decode, six stochastic decodes of the original prompt, one stochastic decode of each of six validated paraphrases, and three additional second decodes that support a hybrid condition. From this single call matrix we derive self-consistency ensembles of any size up to six, paraphrase ensembles of the same sizes, a hybrid allocation, a combined twelve-vote ensemble, per-question consistency scores for both diversity sources, and fragility measures that record how often paraphrasing alone flips the majority answer. Three follow-up conditions then isolate the factors the matrix alone cannot separate: a greedy-paraphrase control that removes decoding stochasticity from the paraphrase side, a chain-of-thought variant that changes the answering mode, and a hard-task probe on MMLU-Pro. Because all conditions are derived from matched calls, comparisons between strategies hold the question, the model, and the total call budget constant by construction.

Our study addresses four research questions. First, under equal budgets, does spending calls on paraphrases buy more, equal, or less accuracy than spending them on resampling, and does the hybrid allocation dominate either pure strategy? Second, how brittle are model answers to paraphrasing per se, and where do paraphrase-induced flips concentrate? Third, can agreement across paraphrases serve as an error signal in the way that agreement across samples does, and how do the two signals compare when both are available? Fourth, how do all of these effects interact with model capability and task difficulty? The last question turns out to govern the first three.

The empirical picture that emerges has a clear regime structure. When a model finds a task easy, its answer distribution collapses: across our six benchmarks, flagship models return the same answer for essentially every sample and every paraphrase on saturated tasks, self-consistency and paraphrase ensembling both collapse onto the greedy baseline, and consistency scores carry no signal about which answers are wrong. When a task genuinely challenges a model, answer entropy reappears, both ensemble families improve over greedy decoding, the two consistency measures flag partially disjoint errors, and a hybrid allocation protects against a failing diversity source. This structure implies that no single published claim about ensembling generalizes across regimes, and that practical systems should first measure the cheap statistics that identify the regime before choosing how to spend their budget.

Our contributions are as follows. We formalize test-time ensembling as a choice between two diversity sources under a shared budget, and we introduce a call matrix design that derives all conditions of interest from one matched set of calls, with controls that separate prompt-space effects from decoding stochasticity, answering mode, and task difficulty. We provide a multi-model, multi-benchmark empirical map of when each source helps, including systematic evidence that paraphrase agreement functions as an uncertainty signal whose error coverage only partially overlaps that of sample agreement. We document the answer degeneracy phenomenon that renders both ensembling and consistency-based uncertainty estimation inert on saturated tasks for confident models, with direct implications for sampling-based confidence estimation in this regime. Finally, we translate the findings into an adaptive allocation rule validated by simulation, and we release our complete prompt materials, validated paraphrases, and per-call predictions.

The remainder of the paper is organized as follows. Section~\ref{sec:related} reviews related work. Section~\ref{sec:method} formalizes the ensembling framework and the protocol. Section~\ref{sec:setup} describes models, benchmarks, and validation procedures. Section~\ref{sec:results} presents the results around our four research questions. Section~\ref{sec:discussion} discusses implications and limitations, and Section~\ref{sec:conclusion} concludes.
